\documentclass[12pt]{article}
\usepackage[utf8]{inputenc}
\usepackage[T1]{fontenc}
\usepackage{hyperref}
\title{Standard Malay Language Resources for K-12 Curriculum Development}
\author{Compiled by Claude Sonnet 5.5 (claude-sonnet-5-5) for Andr\'as Kornai}
\date{}
\begin{document}
\maketitle

\section{Introduction}
Standard Malay (ISO 639-3 zsm), or Bahasa Melayu, is the national language of Malaysia and a language of instruction. Resources are plentiful. My searches covered only the main national institutions: the Dewan Bahasa dan Pustaka (DBP) and the Ministry of Education (KPM). They did not cover the many commercial and community resources.

\section{Primary Educational Resources}
\begin{itemize}
\item KPM digital textbooks on the DELIMa platform. News reports give counts of 754 titles in 2023 and 788 later (352 secondary, 436 primary). \url{https://www.utusan.com.my/berita/2023/03/kpm-hasilkan-70-peratus-buku-teks-digital/}
\item Form 2 Bahasa Malaysia textbook following the DSKP curriculum standard for KSSM, found as a scan. \url{https://anyflip.com/iwzbp/mxrq/basic}
\end{itemize}

\section{Linguistic and Reference Materials}
\begin{itemize}
\item \emph{Kamus Dewan}, DBP's standard dictionary. The first edition appeared in 1970 and the fourth in 2005, with over 49,000 entries. \url{https://www.wikipedia.org/wiki/Kamus_Dewan}
\item \emph{Kamus Dewan Perdana} (2020), DBP's newer dictionary, reported at over 120,000 entries and sub-entries. \url{https://www.utusan.com.my/gaya/sastera/2025/01/kamus-dewan-perdana-dbp-perkukuh-khazanah-tamadun-ilmu/}
\item Overview of DBP's online reference portal PRPM, covering dictionaries, terminology, a thesaurus and spelling checks. \url{https://www.utusan.com.my/nasional/2026/07/usaha-dbp-perkasa-bahasa-melayu-mendepani-kepantasan-dunia-digital/}
\end{itemize}

\section{Religious and Community Resources}
None were examined in this search.

\section{Digital and Audio Resources}
\begin{itemize}
\item DELIMa is KPM's digital learning platform, launched in 2020. KPM also released a free eTextbook Reader app for EPUB textbooks in 2021. \url{https://www.bharian.com.my/amp/berita/nasional/2021/01/775929/bhplus}
\end{itemize}

\section{Curriculum Development Context}
Standard Malay is already taught and used as a medium of instruction in Malaysia under the KSSR and KSSM standards. A curriculum project would mostly build on KPM and DBP materials. Officials report that the digital rollout is limited by device ownership and school internet access.

\end{document}